\SourceLeaf{199}
Insofern die Abbildung gewissen Bedingungen Genüge leisten soll, werden diese Functionen nicht mehr willkührlich sein dürfen. Indem dadurch auch \(X,Y,Z\) zu Functionen von \(t\) und \(u\) werden, müssen diese Functionen, neben der Bedingung, welche die Natur der zweiten Fläche vorschreibt, auch noch derjenigen Genüge leisten, welche in der Abbildung erfüllt werden soll.

\Article{4}
Die Aufgabe der königlichen Gesellschaft der Wissenschaften schreibt vor, dass die Abbildung dem Abgebildeten in den kleinsten Theilen ähnlich sein soll. Es kommt zuvörderst darauf an, diese Bedingung analytisch auszudrücken.

Aus der Differentiation der Functionen von \(t,u\), durch welche \(x,\allowbreak y,\allowbreak z,\allowbreak X,\allowbreak Y,\allowbreak Z\) ausgedrückt werden, mögen folgende Gleichungen hervorgehen:
\[
\begin{aligned}
\mathrm{d}x&=a\,\mathrm{d}t+a'\,\mathrm{d}u\\
\mathrm{d}y&=b\,\mathrm{d}t+b'\,\mathrm{d}u\\
\mathrm{d}z&=c\,\mathrm{d}t+c'\,\mathrm{d}u\\
\mathrm{d}X&=A\,\mathrm{d}t+A'\,\mathrm{d}u\\
\mathrm{d}Y&=B\,\mathrm{d}t+B'\,\mathrm{d}u\\
\mathrm{d}Z&=C\,\mathrm{d}t+C'\,\mathrm{d}u
\end{aligned}
\]
\noindent Die vorgeschriebene Bedingung erfordert, erstlich, dass alle von Einem Punkte der ersten Fläche ausgehende und in ihr liegende unendlich kleine Linien den ihnen entsprechenden Linien der zweiten Fläche proportional sind, und zweitens, dass jene unter sich dieselben Winkel machen, wie diese.

Ein solches Linear-Element auf der ersten Fläche wird
\[
\begin{aligned}
={}&\BareRoot{\bigl((aa+bb+cc)\mathrm{d}t^2+2(aa'+bb'+cc')\mathrm{d}t\mathord{.}\mathrm{d}u}\\
&\qquad+(a'a'+b'b'+c'c')\mathrm{d}u^2\bigr)
\end{aligned}
\]
\noindent und das entsprechende auf der zweiten Fläche
\[
\begin{aligned}
={}&\BareRoot{\bigl((AA+BB+CC)\mathrm{d}t^2+2(AA'+BB'+CC')\mathrm{d}t\mathord{.}\mathrm{d}u}\\
&\qquad+(A'A'+B'B'+C'C')\mathrm{d}u^2\bigr).
\end{aligned}
\]
\noindent Sollen beide, unabhängig von \(\mathrm{d}t\) und \(\mathrm{d}u\), in einem bestimmten Verhältniss zu einander stehen, so müssen offenbar die drei Grössen
\[
aa+bb+cc,\quad aa'+bb'+cc',\quad a'a'+b'b'+c'c'
\]
\noindent respective den drei folgenden proportional sein:
